\Needspace{7\baselineskip}
\pdfbookmark[0]{Anzeigen. Schwab und Metternich: Parallel-Linien. 1816.}{anzeigen-schwab-metternich-1816}
\begin{center}
\rule{.85\textwidth}{.3pt}\par
Göttingische gelehrte Anzeigen.\quad 20. April 1816.\par
\rule{.85\textwidth}{.3pt}
\end{center}
\hspace*{\parindent}\textit{Commentatio in primum elementorum} \textsc{Euclidis} \textit{librum, qua veritatem geometriae principiis ontologicis niti evincitur, omnesque propositiones, axiomatum geometricorum loco habitae, demonstrantur. Auctore} \textsc{J. C. Schwab}, \textit{Regi Württembergiae a consiliis aulicis secretioribus, academiae scientiarum Petropolitanae, Berolinensis et Harlemensis Sodali.} (65 Seiten im Oktavformat.) Stuttgart 1814. Typis \textsc{J. F. Steinkopf}.

\textit{Vollständige Theorie der Parallel-Linien. Nebst einem Anhange, in welchem der erste Grundsatz zur Technik der geraden Linie angegeben wird. Herausgegeben von} \textsc{Matthias Metternich}, \textit{Doctor der Philosophie, Professor der Mathematik, Mitglied der gelehrten Gesellschaft nützlicher Wissenschaften zu Erfurt.} 44 Seiten im Oktavformat. Mainz 1815. Auf Kosten des Verfassers in Kommission bei \textsc{Florian Kupferberg}.

Über wenige Gegenstände der Mathematik dürfte so viel geschrieben worden sein wie über die Lücke am Anfang der Geometrie bei der Begründung der Theorie der Parallelen. Selten vergeht ein Jahr, in dem kein neuer Versuch unternommen wird, diese Lücke zu schließen. Wenn wir ehrlich und offen sprechen wollen, können wir aber nicht sagen, dass wir im Wesentlichen weitergekommen wären als \textsc{Euklid} vor 2000 Jahren. %
